\chapter{Les cours de Buffalo, 1973}\label{ch:buffalo}

\section{La source}

Au printemps et à l'été 1973, Grothendieck a donné trois cours à l'université d'État de
New York à Buffalo : la théorie des topos (d'avril à juillet), la géométrie algébrique (de
mai à juillet) et les groupes algébriques (en juin et juillet). Les cours sont en anglais.
Des bandes magnétiques les ont enregistrés ; John Duskin et William Lawvere les ont
gardées, et une chaîne de vidéos les a publiées en octobre 2025, en dix vidéos (quatre pour
les topos, quatre pour la géométrie algébrique, deux pour les groupes), soit environ
89~heures. Chaque vidéo enchaîne plusieurs faces de bande ; les descriptions donnent la
date de chaque séance, le numéro de la bande, la face, et signalent les faces abîmées.

C'est la seule source de ce livre qui soit du son, et la seule dont l'image ne montre pas
ce qui compte : le cours se faisait au tableau noir, et les bandes n'ont gardé que la voix.

\section{Les particularités d'une source orale}

\paragraph{Pas de page à transcrire.} Il n'y a pas de fac-similé à afficher à côté d'un
texte. Le site joue la vidéo depuis YouTube, et aucun son n'est copié. Un horodatage du
site ou d'un PDF est un lien vers la vidéo, à la seconde.

\paragraph{La machine entend, elle ne voit pas.} Les sous-titres automatiques de la plupart
des vidéos ont servi de première lecture ; la deuxième vidéo de topos (du 30~mai à la fin de
juin) est trop bruitée pour cela, et a été lue sur une transcription complète de la piste
sonore, faite en local par reconnaissance vocale. Ces transcriptions restent dans un
dossier ignoré par \texttt{git}. Elles se trompent souvent, surtout sur les noms propres et
les formules dictées. Les entrées ne sont donc pas toutes réécoutées une à une : la
relecture contre la bande est faite par sondage, et le site ne la déclare pas plus large
qu'elle n'est.

\paragraph{Les zones d'ombre.} Les voix ne sont pas nommées. Colin McLarty, qui n'y était
pas mais connaît ces bandes, en identifie quelques-unes, et le site dit ce qui
vient de lui et ce qui est une déduction du projet. Les dates des séances de la deuxième
vidéo se lisent dans le texte et non dans la description. Une seule date, la dernière,
fait question : les descriptions donnent le 12~juillet 1973 pour toute la quatrième vidéo
de topos, McLarty le 13. Le projet suit les descriptions et le dit. Quatre faces de bande
sont trop dégradées pour être indexées.

\paragraph{Un cours qu'on ne cite pas.} À l'oral on dit des choses approximatives, parfois fausses,
et l'ordre de l'exposé suit la salle plus que la logique. Une transcription fidèle
reproduirait tout cela ; le projet a préféré ne pas transcrire.

\section{Les droits, et ce qu'on a décidé}

Un cours est une œuvre orale ; les droits sont ceux des ayants droit de Grothendieck, et
la réponse de la BnF sur le fonds de Lasserre (aucune transcription en ligne sans leur
autorisation) a été appliquée ici par prudence. Le projet ne publie donc :
\begin{itemize}
  \item aucune transcription du cours, ni aucun extrait de la piste sonore ;
  \item des énoncés mathématiques \emph{standards, écrits par nous}, un par ligne, avec le
        moment de la bande, et quelques mots entre guillemets quand le cours a forgé un
        nom ;
  \item l'énoncé correct, avec ses hypothèses, quand ce qu'on entend est une esquisse, et
        jamais une affirmation fausse telle quelle ;
  \item des notes de tableau, « ce qui a probablement été écrit », dans nos mots, présentées
        comme une reconstitution à partir des sous-titres et non comme le tableau.
\end{itemize}
Les enregistrements, eux, sont ceux d'un tiers : on renvoie vers la vidéo et on ne la copie
pas.

\section{Le bilan}

Au début de novembre 2026 :
\begin{itemize}
  \item trois index de résultats, un par cours, en LaTeX : 807 entrées pour les topos, 982
        pour la géométrie algébrique et 335 pour les groupes, en 225 rubriques. Chaque
        entrée est un énoncé (définition, proposition, théorème, exemple,
        contre-exemple, remarque, question) avec son horodatage ;
  \item soixante-treize entrées « Discussion » : des échanges avec la salle où un point est
        tranché, corrigé ou laissé ouvert, décrits sans citer ;
  \item trente entrées « Lecture » aux moments où il recommande un livre, renvoyant à une
        bibliographie numérotée de vingt-quatre titres (Godement, les SGA, Giraud, Demazure
        et Gabriel, Mumford…), entre crochets comme dans un article ;
  \item pour chaque cours, un PDF où l'horodatage de chaque entrée est en marge, à droite,
        petit et pâle, et où les théorèmes sont encadrés ; il est compilé depuis le même
        fichier que la page du site, si bien que les deux ne peuvent pas diverger ;
  \item une page du site, avec le lecteur vidéo, l'index à côté, et un tableau noir
        reconstitué (121 tableaux) dessiné comme celui d'une salle de l'IHP.
\end{itemize}

Le cours de topos a été relu contre la conférence de McLarty de 2018 : plusieurs
formules des notes et du tableau en ont été corrigées, et sa remarque sur les voix et la
date de la dernière séance est dans l'introduction de l'index. Une vérification contre ses
diapositives reste à faire ; la lecture des images d'une vidéo demande un téléchargement
que le projet n'a pas encore décidé de faire.

\section{Les éléments qu'on peut en extraire}

Un index n'est pas la seule chose qu'une telle source permet. Les éléments qui reviennent
d'un cours à l'autre, et qu'on peut extraire d'une bande sans la transcrire, sont les
suivants.
\begin{itemize}
  \item \textbf{Les énoncés}, avec leur nature (définition, théorème, exemple, remarque,
        question) et le moment où ils sont traités. C'est la matière des trois index.
  \item \textbf{Les corrections.} Quand il esquisse ou se trompe, l'énoncé juste avec ses
        hypothèses, et une note quand l'énoncé entendu ne tient pas.
  \item \textbf{Les lectures recommandées}, avec l'instant où il les dit : une bibliographie
        du cours, que personne n'avait dressée.
  \item \textbf{Les discussions avec la salle}, où un point se règle ou reste ouvert ; les
        questions posées disent aussi ce que le public savait.
  \item \textbf{Les notions laissées floues}, comme l'« hypertopos » ou le « pseudo-topos » :
        on les signale et on n'en fait pas une définition.
  \item \textbf{Le plan}, enfin : l'ordre des séances, les dates, les faces perdues. Il
        permet de situer un énoncé dans le cours et de le relier à d'autres fonds (les
        dossiers de Montpellier qui y ressemblent, les rédactions Bourbaki).
\end{itemize}
Chacun de ces éléments se vérifie à la seconde près en rouvrant la vidéo, ce qu'aucune
transcription de brouillon ne permet aussi simplement.

\begin{exemple}
Pour une source orale protégée, l'index est le bon format. La relecture y demande surtout
de savoir les mathématiques, là où un manuscrit demande de savoir lire une main.
\end{exemple}
